\section{Introduction}

The congruence closure data structure,
 also known as the \egraph,
 is a central component of SMT-solvers~\cite{simplify, z3, moskal, cvc4}
 and equality saturation-based optimizers~\cite{eqsat, egg}.
An \egraph compactly represents a set of terms
 and an equivalence relation over the terms.
An important operation on \egraphs is \ematching,
 which finds the set of terms in an \egraph
 matching a given pattern.
In SMT-solvers, \ematching is used to instantiate
 quantified formulas over ground terms.
% For instance, suppose the solver input contains the
%  formula $\forall x . \phi(x, g(x)) \rightarrow \psi(x)$
%  for predicate symbol $\phi, \psi$
%  and function symbol $g$.
% The solver will search for potential instantiation
%  of the form $\phi(x, g(x))$ in its internal \egraph
%  data structure.
% For every ground term found,
%  the solver is able to deduce $\psi(x)$ as a fact.
In equality saturation, \ematching is used
 to match rewrite rules on an \egraph
 to discover new equivalent programs.
% In equality saturation, rewrite rules of the form
%  $f(\alpha, g(\alpha)) \Rightarrow h(\alpha)$ match
%  the left-hand-side pattern $f(\alpha, g(\alpha))$
%  on terms represented by the \egraph;
% upon every match, the right-hand-side $h(\alpha)$ is
%  instantiated with the substitution $\alpha \mapsto t$
%  into $h(t)$ and is added to the \egraph.
% One can view the rewrite rule
%  $f(\alpha, g(\alpha)) \Rightarrow h(\alpha)$
%  as a formula where the pattern variable $\alpha$ is
%  $\forall$-quantified,
%  so the \ematching~instances in SMT-solvers and equality
%  saturation are practically the same.
The efficiency of \ematching~greatly affects the overall performance
 of the SMT-solver~\cite{z3, cvc4},
 and slow \ematching~is a major bottleneck
 in equality saturation~\cite{egg, tensat, szalinski}.
In a typical application of equality saturation,
 \ematching~is responsible for 60--90\% of the overall run time~\cite{egg}.

Several algorithms have been proposed for e-matching~\cite{efficient-ematching, moskal, simplify}.
However, due to the NP-completeness of \ematching~\cite{ematching-nph},
 most algorithms implement some form of backtracking search,
 which are inefficient in many cases.
In particular, backtracking search only exploits \textit{structural constraints},
 which are constraints about the shape of a pattern,
 but defers checking \textit{equality constraints},
 which are constraints that variables should be consistently mapped.
This leads to suboptimal run time when
 the equality constraints dominate the structural constraints.

To improve the performance of backtracking-based \ematching,
 existing systems implement various optimizations.
Some of these optimizations only deal with
 patterns of certain simple shapes and are therefore \textit{ad hoc}
 in nature~\cite{moskal}.
Others attempt to incrementalize \ematching\ upon changes to the \egraph,
 or match multiple similar patterns together
 to eliminate duplicated work~\cite{efficient-ematching}.
However, these optimizations are complex to implement
 and fail to generalize to workloads where the \egraph~changes rapidly
 or when the patterns are complex.

% Finally, despite the NP-completeness of the general \ematching~problem, most
% researchers acknowledge \ematching~is almost always tractable, thanks to the
% small size of patterns encountered in practice. This hints at possible
% algorithmic improvements to \ematching.

To tackle the inefficiency and complexity involved in \ematching,
 we propose a systematic approach to \ematching~called \textbf{relational \ematching}.
Our approach is based on the observation that
 \ematching\ is an instance of a well-studied problem
 in the databases community,
 namely answering {\em conjunctive queries}.
We therefore propose to solve \ematching\ on an \egraph\ by reducing it
 to answering conjunctive queries on a relational database.
This approach has several benefits.
First, by reducing \ematching to conjunctive queries,
 we simplify \ematching by taking advantage of
 decades of study by the databases community.
Second, the relational representation provides a unified way to
 express both the structural constraints and equality constraints in patterns,
 allowing query optimizers to leverage both kinds of constraints
 to generate asymptotically faster query plans.
Finally, by leveraging the generic join algorithm,
 a novel algorithm developed in the databases community,
 our technique achieves the first worst-case optimal bound for \ematching.

Relational \ematching is provably optimal despite the NP-hardness of \ematching.
The databases community makes a clear distinction between
 {\em query complexity}, the complexity dependent on the size of the query,
 and {\em data complexity}, the complexity dependent on the size of the database.
The NP-hardness result~\cite{ematching-nph} is stated over the size of the pattern,
 yet in practice only small patterns are matched on a large \egraph.
When we hold the size of each pattern constant,
 relational \ematching runs in time polynomial over the size of the \egraph.

Our approach is widely applicable.
For example, \textit{multi-patterns} are
 typically framed as an extension to \ematching
 that allows the user to find matches
 satisfying multiple patterns simultaneously.
Efficient support for multi-patterns requires modifying the basic
 backtracking algorithm~\cite{efficient-ematching}.
In contrast, relational \ematching inherently supports multi-patterns for free.
The relational model also opens the door
 to entirely new kinds of optimizations,
 such as persistent or incremental \egraphs.

To evaluate our approach, we implemented relational \ematching for \egg, a
 state-of-the-art implementation of \egraphs.
Relational \ematching
 is simpler, more modular, and orders of magnitude faster
 than \egg's \ematching implementation.

In summary, we make the following contributions in this paper:
\begin{itemize}
\item We propose relational \ematching, a systematic approach to \ematching that
  is simple, fast, and optimal.
\item We adapt generic join to implement relational \ematching,
  and provide the first data complexity results for \ematching.
\item We prototyped relational \ematching\footnote{
  Our implementation can be found here: \url{https://github.com/yihozhang/relational-ematching-benchmark}.
  } in \egg,
  a state-of-the-art \egraph implementation,
  and we show that relational \ematching can be orders of magnitude faster.
\end{itemize}

The rest of the paper is organized as follows: \autoref{sec:background} reviews
relevant background on the \egraph~data structure, the \ematching~problem,
conjunctive queries and join algorithms.
\autoref{sec:algorithm} presents our
 relational view of \egraphs,
 our \ematching algorithm,
 and the complexity results.
\autoref{sec:optimization} discusses
optimizations on our core algorithm and addresses various practical concerns.
\autoref{sec:evaluation} evaluates our algorithm and implementation with a set of
experiments in the context of equality saturation.
\autoref{sec:discussion} discusses how the relational model opens up
 many avenues for future work in \egraphs and \ematching,
 and \autoref{sec:conclusion} concludes.

%%% Local Variables:
%%% mode: latex
%%% TeX-master: "main"
%%% End:
